\documentclass{article}

\usepackage[final]{style}
\usepackage[utf8]{inputenc}
\usepackage[T1]{fontenc}
\usepackage{hyperref}
\usepackage{url}
\usepackage{booktabs}
\usepackage{tabularx}
\usepackage{amsfonts}
\usepackage{amsmath}
\usepackage{nicefrac}
\usepackage{microtype}
\hypersetup{hidelinks}

% The provided course style leaves these footer fields undefined.
\makeatletter
\providecommand{\@neuripsordinal}{Fall 2026}
\providecommand{\@neuripslocation}{Carnegie Mellon University}
\makeatother

\title{AirClay: Few-Shot, Depth-Aware Gesture Control\\
for 3D Sculpting in Blender}

\author{%
  Author 1 \\
  Carnegie Mellon University \\
  \texttt{author1@andrew.cmu.edu} \\
  \And
  Author 2 \\
  Carnegie Mellon University \\
  \texttt{author2@andrew.cmu.edu} \\
  \And
  Author 3 \\
  Carnegie Mellon University \\
  \texttt{author3@andrew.cmu.edu} \\
  \AND
  Author 4 \\
  Carnegie Mellon University \\
  \texttt{author4@andrew.cmu.edu} \\
  \And
  Author 5 \\
  Carnegie Mellon University \\
  \texttt{author5@andrew.cmu.edu} \\
}

\begin{document}

\maketitle

\begin{abstract}
We propose AirClay, a markerless two-hand interface for sculpting a deformable
mesh in Blender with a commodity RGB webcam. The system will solve two coupled
vision problems: causal recognition of four continuous interaction intents
(\textit{idle}, \textit{sculpt}, \textit{orbit}, and \textit{resize}) and
estimation of the hand's short-term displacement in camera space. Public
gesture and 3D hand-pose datasets will provide most of the training data, while
one, five, or ten demonstrations per AirClay command will personalize a
prototype-based temporal classifier. We will compare rules, a single-frame
MLP, a causal GRU trained from scratch, public-data pretraining, and few-shot
adaptation. We will separately compare geometric, hand-specific learned, and
generic monocular depth cues. Evaluation will measure recognition quality,
false activations, 3D displacement error, latency, and performance on a fixed
Blender sculpting task. Our main question is whether public-data pretraining
and few-shot personalization can enable reliable, depth-aware two-hand
sculpting while preserving interactive response time.
\end{abstract}

\section{Problem Statement}

Three-dimensional content is increasingly central to animation and games,
industrial and product design, digital fabrication, scientific visualization,
and immersive AR/VR experiences. Yet producing that content remains a skilled
task. In conventional desktop software, a user controls a 3D surface through a
2D mouse or tablet while repeatedly coordinating camera navigation, brush
parameters, keyboard modifiers, and editing modes. These tools provide high
precision, but their indirect mapping can make rapid shape exploration less
natural: the creator's intended hand motion must first be translated into a
sequence of device movements and mode changes.

Bare-hand interaction offers a complementary way to specify shape. It can map
the motion of the fingers and palms directly to the location, direction, and
magnitude of an edit; it also permits asymmetric bimanual interaction, where
one hand edits while the other controls the view or scale. Using an ordinary
RGB webcam would avoid instrumented gloves, controllers, and dedicated depth or
motion-capture hardware, lowering the equipment and setup requirements for
casual modeling and classroom use. These are potential interaction benefits,
however, and they depend on sufficiently reliable perception. A small number
of visually intuitive commands is only useful if natural hand motion does not
silently alter the model.

Existing approaches leave several gaps. Gesture-controlled deformation systems
often simplify perception with colored gloves or other markers
\citep{mahdikhanlou2022deforming}. Clay-motion recognition has been studied
with accurate multi-camera Vicon trajectories \citep{veldhuijzen2024shrec}, but
such inputs do not represent the noise, occlusion, and missing detections of a
consumer webcam. Many gesture benchmarks classify pre-segmented clips or a
fixed closed set, whereas an editor receives an unsegmented stream containing
transitions and unknown actions. Their clip-level accuracy therefore does not
measure the false activations, unintended releases, or onset delay that directly
affect an editing session. Finally, monocular hand-pose methods can recover
articulation, but absolute scale and camera-space depth remain ambiguous from a
single RGB view, especially when two hands overlap or move toward the camera
\citep{huang2023nvf,yu2023acr}.

AirClay will address these gaps as a causal, depth-aware control problem. Given
RGB frames observed up to time $t$, the system will estimate an intent
$y_t \in \{\textit{idle},\textit{sculpt},\textit{orbit},\textit{resize}\}$,
a short-term camera-space displacement
$\Delta\mathbf{p}_t=(\Delta x_t,\Delta y_t,\Delta z_t)$, and a confidence value.
The dominant hand will pinch and drag to deform a local surface region, the
non-dominant hand will orbit the model, and a two-hand pinch-and-spread motion
will resize the brush. Low-confidence predictions and sustained hand loss will
pause or roll back an operation instead of modifying the mesh.

We will test two linked hypotheses. First, a causal temporal encoder pretrained
on public gesture data and adapted with one, five, or ten user demonstrations
will reduce false activations and premature releases relative to thresholds, a
single-frame MLP, and a GRU trained only on the small target set. Second, a
calibrated hand-specific depth estimator will recover forward/backward motion
more accurately and with less jitter than planar control, inverse palm scale,
or a generic monocular depth model. The downstream test will require a user to
create two protrusions, rotate the object, resize the brush, undo one edit, and
save the result in Blender. Thus the project will connect recognition and depth
errors to task completion time, edit accuracy, undo count, and interaction
latency rather than evaluating perception only in isolation.

\subsection{Related Work}

\paragraph{Continuous and dynamic gesture recognition.}
Jester contains 148,092 short RGB clips of 27 gestures collected in front of
laptop or webcam cameras, including explicit non-gesture classes
\citep{materzynska2019jester}. IPN Hand instead targets continuous recognition:
its 50 subjects perform 4,218 gestures in more than 800,000 frames with natural
motion between commands \citep{benitez2021ipn}. The recent IPN HandS extension
adds refined temporal boundaries and frame-level 21-joint skeleton annotations
\citep{benitez2025ipnhands}. These datasets make it possible to learn temporal
features and rejection behavior before collecting application-specific data.
HaGRID provides complementary diversity in hand appearance, background,
illumination, and camera distance across tens of thousands of people
\citep{kapitanov2024hagrid}. Recent multimodal attention work further shows the
value of combining skeleton and appearance features for dynamic gestures
\citep{balaji2024multimodal}. However, most benchmarks report clip accuracy;
AirClay will also evaluate false activations, release errors, and latency in an
unsegmented interactive stream.

\paragraph{Clay motion and direct 3D manipulation.}
The SHREC 2024 clay-motion benchmark is the closest domain-specific work. It
contains 62 motions in seven fine-grained pottery classes, recorded from one
expert potter using 14 Vicon cameras and 28 markers across both hands
\citep{veldhuijzen2024shrec}. Its methods include recurrent networks, HMMs,
spatiotemporal graph networks, and a skeleton masked autoencoder. This study
demonstrates that clay motions are long, precise, and visually similar, but it
classifies isolated motion-capture sequences and does not control an editor.
Mahdikhanlou and Ebrahimnezhad deform a 3D object with dynamic gestures from an
RGB camera, but use colored gloves to simplify hand observation
\citep{mahdikhanlou2022deforming}. AirClay will instead use bare hands,
continuous causal recognition, and real Blender mesh operations.

\paragraph{Few-shot and user-defined gestures.}
Few-shot learning is attractive because a fixed public taxonomy cannot encode
every editor command or a user's preferred motion. Liu et al. combine
MediaPipe landmarks, lightweight autoencoders, attention, and similarity-based
classification for few-shot continuous RGB gestures
\citep{liu2022continuous}. Schluesener and Buecker train LSTM Relation Networks
episodically on landmark sequences derived from Jester, then recognize new
gesture sets from one, two, or five examples \citep{schluesener2022fast}.
Zhou et al. address open-set hand gestures with viewpoint-independent features
and cosine-based few-shot incremental classification
\citep{zhou2025openset}. Building on these ideas, AirClay will learn a general
motion embedding from public datasets and adapt only a small projection head
and class prototypes from user demonstrations.

\paragraph{Monocular depth and interacting hands.}
MediaPipe provides useful hand-centered 3D landmarks, but each hand's local
coordinates do not determine a common camera-space translation. Absolute depth
from one RGB view is inherently ambiguous. MobRecon demonstrates fast
single-view hand-mesh recovery with temporal coherence
\citep{chen2022mobrecon}, while Neural Voting Field directly estimates
camera-space 3D joints and global hand geometry \citep{huang2023nvf}. For two
hands, ACR explicitly handles mutual occlusion, truncation, and cross-hand
confusion \citep{yu2023acr}. Generic monocular models such as Depth Anything V2
offer robust relative scene depth but are not designed for millimeter-scale
hand motion \citep{yang2024depthanything}. We will therefore compare a generic
depth model with a geometric palm-scale baseline and a compact hand-specific
depth regressor, and will focus on calibrated short-term displacement rather
than claiming uncalibrated metric depth.

\section{Data}

Public data will supply broad appearance, motion, and depth supervision. A
small AirClay adaptation set will be used only for editor-specific semantics
and per-user scale calibration. RHD supplies paired RGB, camera calibration,
depth, and 3D joints \citep{zimmermann2017hand}, while InterHand2.6M supplies
camera-space annotations for interacting hands \citep{moon2020interhand}.
Table~\ref{tab:data} summarizes the planned sources.

\begin{table}[h]
\centering
\footnotesize
\setlength{\tabcolsep}{4pt}
\begin{tabularx}{\textwidth}{@{}p{0.17\textwidth}p{0.35\textwidth}X@{}}
\toprule
Dataset & Size and annotation & Planned use \\
\midrule
\href{https://www.qualcomm.com/developer/software/jester-dataset}{Jester} & 148,092 RGB clips, 27 classes, 22.8 GB & Pretrain dynamic motion and gesture rejection from webcam-like video. \\
\href{https://gibranbenitez.github.io/IPN_Hand/}{IPN Hand / HandS} & 50 subjects, 4,218 instances, $>$800K continuous RGB frames, boundaries, 21-joint skeletons & Pretrain causal spotting, landmark sequences, and idle/hard-negative behavior. \\
\href{https://github.com/o-ikne/skelmae-shrec24}{SHREC 2024 Clay} & 62 sequences, 7 pottery classes, bilateral 28-marker 3D tracks & Domain-specific contrastive pretraining on subtle two-hand motion. \\
\href{https://lmb.informatik.uni-freiburg.de/resources/datasets/RenderedHandposeDataset.en.html}{RHD} & 41,258 train and 2,728 test samples; RGB, depth, intrinsics, segmentation, and 3D joints; 7.1 GB & Train and test an RGB/landmark-to-camera-depth regressor. \\
\href{https://mks0601.github.io/InterHand2.6M/}{InterHand2.6M} & 2.59M single- and interacting-hand RGB frames with camera and 3D joint annotations & Use a fixed sample from the official 5 FPS interacting-hand subset for two-hand depth and occlusion tests. \\
\href{https://github.com/hukenovs/hagrid}{HaGRID} & Large RGB gesture collection across subjects, scenes, and 0.5--4 m distance & Use public weights or a fixed class-balanced subset for appearance robustness. \\
\bottomrule
\end{tabularx}
\caption{Public datasets and their distinct roles. Large optional datasets will
be sampled with stored frame IDs so that experiments remain reproducible.}
\label{tab:data}
\end{table}

The core reproducible pipeline will require Jester, IPN HandS, SHREC 2024, and
RHD. We will not require the full 1.5 TB HaGRIDv2 release; we will use public
weights or selected classes. Likewise, InterHand2.6M will be restricted to a
fixed, subject- and view-stratified sample of its 5 FPS interacting-hand data.
Every RGB dataset will be processed by one pinned MediaPipe version
\citep{zhang2020mediapipe} into a
common representation with up to 21 image landmarks per hand, timestamps,
handedness, and validity masks. When curated skeletons are available, both the
MediaPipe predictions and ground truth will be retained. Dataset-native 3D
coordinates will remain in separate fields with their camera calibration;
incompatible coordinate systems will not be concatenated as if they shared a
scale.

Public labels will supervise reusable primitives---idle, pinch/click,
translation, push/pull, and two-hand scale---rather than being relabeled
directly as Blender commands. For each target user, we will then record one,
five, or ten demonstrations of each AirClay intent plus two minutes of natural
hand motion. Three short calibration sequences at approximately 30, 45, and
60 cm will estimate person- and camera-specific scale. Support, validation, and
query sequences will be separated before sliding windows are generated, and
the final test will hold out people rather than frames.

\section{Method}

\subsection{Representation and Public-Data Pretraining}

At roughly 30 Hz, the input will include wrist-centered, palm-scale-normalized
landmarks; wrist and fingertip image positions and velocities; pinch ratios;
two-hand distance and relative velocity; hand-presence masks; and lightweight
features from hand RGB crops. All online features will use only the current and
previous frames. A causal GRU or temporal convolutional network (TCN) will
encode a 24-frame window into a motion embedding. We will first train this
encoder on the Jester and IPN primitives with a supervised classification loss.
SHREC sequences will add a contrastive objective so that clips from the same
fine-grained clay class remain close despite duration changes. Public
non-gesture segments and failed detections will supply rejection examples.

\subsection{Few-Shot Intent Recognition}

Our primary few-shot model will use a Prototypical Network. If
$f_{\theta}(x)$ is the temporal embedding for sequence $x$, the prototype for
class $c$ will be
\begin{equation}
  \mathbf{p}_c = \frac{1}{|S_c|}\sum_{x_i \in S_c} f_{\theta}(x_i),
\end{equation}
where $S_c$ is the user's support set. Prediction will use cosine distance to
the four prototypes. A validation-set distance threshold will reject
low-confidence windows as \textit{idle}. During personalization, most of the
public-data encoder will remain frozen; only the projection layer and class
prototypes will be updated. Episodic training on public gesture classes will
make the training objective match 1-, 5-, and 10-shot evaluation.

We will compare five recognition approaches: pinch thresholds and hand-role
rules; a single-frame MLP; a causal GRU trained from scratch on the same support
data; the public-pretrained GRU with a linear classifier; and the pretrained
GRU with prototype adaptation. All approaches will feed the same causal state
machine. The state machine will use probability smoothing, onset/offset
hysteresis, a short two-hand arbitration window, mode locking, and rollback
after sustained hand loss. This isolates the recognition model as the
experimental variable.

\subsection{Depth and Continuous Motion}

Intent classification will be separated from continuous control. The latter
will estimate short-term camera-space displacement
\begin{equation}
  \Delta\mathbf{p}_t = (\Delta x_t,\Delta y_t,\Delta z_t).
\end{equation}
Image-plane landmark motion and the calibrated camera projection will provide
$\Delta x_t$ and $\Delta y_t$. We will compare three estimates of
$\Delta z_t$: (1) inverse apparent palm width after three-distance calibration;
(2) a compact regressor trained on RHD and sampled InterHand RGB, 2D joints,
and camera-space root depth; and (3) relative depth sampled from Depth Anything
V2 Small inside the hand region. The learned regressor will predict both depth
and confidence. Low-confidence frames will retain planar motion but suppress
depth updates. Temporal filtering and maximum-velocity limits will prevent a
single unstable estimate from causing a large mesh edit.

Blender will receive intent, normalized pointer position, displacement,
confidence, sequence number, and operation ID over local UDP. In sculpt mode,
a view ray will select the surface, and the camera-space displacement will be
mapped to the current view basis. Nearby vertices will move with a smooth
geodesic falloff from a snapshot taken at gesture onset. Orbit mode will map
motion to camera azimuth and elevation, while resize mode will map the change
in bilateral 3D distance to brush radius. Completed operations will enter a
30-step undo/redo history.

\subsection{Evaluation}

For recognition, we will report macro-F1, per-class precision and recall,
unknown-motion rejection, false activations per minute, missed operations,
premature releases, mode-switch errors, and onset/offset latency. Few-shot
results will be reported at 1, 5, and 10 shots on held-out users. For depth, we
will report wrist/root depth mean absolute error, 3D joint MPJPE where ground
truth is available, forward/backward direction accuracy, temporal jitter, and
inference time.

Finally, participants will perform a fixed Blender task: create two
protrusions on a sphere, rotate the model, resize the brush, correct one edit
with undo, and save the result. We will compare planar-only, geometric-depth,
and learned-depth control using task completion time, success rate, undo count,
overshoot count, and final control-point error. The target is at least 20 Hz
interaction updates and median end-to-end latency below 250 ms; these are
performance goals to be measured, not assumed results.

\begin{thebibliography}{99}
\small

\bibitem[Balaji and Prusty(2024)]{balaji2024multimodal}
P. Balaji and M. R. Prusty.
\newblock Multimodal fusion hierarchical self-attention network for dynamic hand gesture recognition.
\newblock \emph{Journal of Visual Communication and Image Representation}, 98:104019, 2024.

\bibitem[Benitez-Garcia et al.(2021)]{benitez2021ipn}
G. Benitez-Garcia, J. Olivares-Mercado, G. Sanchez-Perez, and K. Yanai.
\newblock IPN Hand: A video dataset and benchmark for real-time continuous hand gesture recognition.
\newblock In \emph{Proceedings of ICPR}, pages 4340--4347, 2021.

\bibitem[Benitez-Garcia et al.(2025)]{benitez2025ipnhands}
G. Benitez-Garcia, J. Olivares-Mercado, G. Sanchez-Perez, and H. Takahashi.
\newblock IPN HandS: Efficient annotation tool and dataset for skeleton-based hand gesture recognition.
\newblock \emph{Applied Sciences}, 15(11):6321, 2025.

\bibitem[Chen et al.(2022)]{chen2022mobrecon}
X. Chen, Y. Liu, Y. Dong, X. Zhang, C. Ma, Y. Xiong, Y. Zhang, and X. Guo.
\newblock MobRecon: Mobile-friendly hand mesh reconstruction from monocular image.
\newblock In \emph{Proceedings of CVPR}, pages 20544--20554, 2022.

\bibitem[Huang et al.(2023)]{huang2023nvf}
L. Huang, C.-C. Lin, K. Lin, L. Liang, L. Wang, J. Yuan, and Z. Liu.
\newblock Neural Voting Field for camera-space 3D hand pose estimation.
\newblock In \emph{Proceedings of CVPR}, pages 8969--8978, 2023.

\bibitem[Kapitanov et al.(2024)]{kapitanov2024hagrid}
A. Kapitanov, K. Kvanchiani, A. Nagaev, R. Kraynov, and A. Makhliarchuk.
\newblock HaGRID: HAnd Gesture Recognition Image Dataset.
\newblock In \emph{Proceedings of WACV}, pages 4572--4581, 2024.

\bibitem[Liu et al.(2022)]{liu2022continuous}
Z. Liu, C. Pan, and H. Wang.
\newblock Continuous gesture sequences recognition based on few-shot learning.
\newblock \emph{International Journal of Aerospace Engineering}, 2022:7868142, 2022.

\bibitem[Mahdikhanlou and Ebrahimnezhad(2022)]{mahdikhanlou2022deforming}
K. Mahdikhanlou and H. Ebrahimnezhad.
\newblock 3D object deforming and manipulating through dynamic hand gestures.
\newblock \emph{Entertainment Computing}, 43:100507, 2022.

\bibitem[Materzynska et al.(2019)]{materzynska2019jester}
J. Materzynska, G. Berger, I. Bax, and R. Memisevic.
\newblock The Jester dataset: A large-scale video dataset of human gestures.
\newblock In \emph{Proceedings of ICCV Workshops}, 2019.

\bibitem[Moon et al.(2020)]{moon2020interhand}
G. Moon, S.-I. Yu, H. Wen, T. Shiratori, and K. M. Lee.
\newblock InterHand2.6M: A dataset and baseline for 3D interacting hand pose estimation from a single RGB image.
\newblock In \emph{Proceedings of ECCV}, pages 548--564, 2020.

\bibitem[Schluesener and Buecker(2022)]{schluesener2022fast}
N. Schluesener and M. Buecker.
\newblock Fast learning of dynamic hand gesture recognition with few-shot learning models.
\newblock arXiv:2212.08363, 2022.

\bibitem[Veldhuijzen et al.(2024)]{veldhuijzen2024shrec}
B. Veldhuijzen, R. C. Veltkamp, O. Ikne, et al.
\newblock SHREC 2024: Recognition of dynamic hand motions molding clay.
\newblock \emph{Computers \& Graphics}, 123:104012, 2024.

\bibitem[Yang et al.(2024)]{yang2024depthanything}
L. Yang, B. Kang, Z. Huang, Z. Zhao, X. Xu, J. Feng, and H. Zhao.
\newblock Depth Anything V2.
\newblock In \emph{Advances in Neural Information Processing Systems}, 37, 2024.

\bibitem[Yu et al.(2023)]{yu2023acr}
Z. Yu, S. Huang, C. Fang, T. P. Breckon, and J. Wang.
\newblock ACR: Attention collaboration-based regressor for arbitrary two-hand reconstruction.
\newblock In \emph{Proceedings of CVPR}, pages 12955--12964, 2023.

\bibitem[Zhang et al.(2020)]{zhang2020mediapipe}
F. Zhang, V. Bazarevsky, A. Vakunov, A. Tkachenka, G. Sung, C.-L. Chang, and M. Grundmann.
\newblock MediaPipe Hands: On-device real-time hand tracking.
\newblock arXiv:2006.10214, 2020.

\bibitem[Zhou et al.(2025)]{zhou2025openset}
J. Zhou, C. Xu, and L. Cheng.
\newblock Hand gesture recognition from an open-set perspective.
\newblock \emph{IEEE Transactions on Multimedia}, 27:4181--4192, 2025.

\bibitem[Zimmermann and Brox(2017)]{zimmermann2017hand}
C. Zimmermann and T. Brox.
\newblock Learning to estimate 3D hand pose from single RGB images.
\newblock In \emph{Proceedings of ICCV}, pages 4903--4911, 2017.

\end{thebibliography}

\end{document}
